\section{External evidence base: how the findings sit in the literature}
\label{sec:evidence}

The preceding chapters establish the FACE-ATLAS results internally. This chapter asks the complementary
question: do they agree with what is already known? The purpose is to show the findings are neither isolated
nor artefactual --- and the stance is \emph{calibration, not advocacy}. Where the published literature
qualifies or contradicts a claim, we say so plainly, because a reviewer will probe exactly those points and the
result is more persuasive for surviving the tension than for ignoring it. Every study cited here was retrieved
and verified against PubMed (title, authors, journal, year, DOI) before inclusion; the full annotated set,
including the papers we checked and set aside, is recorded in \file{docs/LITERATURE_EVIDENCE.md}.

\begin{table}[h]\centering\small\sffamily
\caption{Claim-by-claim verdict against the external literature. ``Supported'' = a substantial independent
literature agrees; ``open'' = an active, unsettled debate the work contributes to; ``calibrated'' = supported
in direction but with a boundary the wording must respect.}
\label{tab:evidence}
\begin{tabular}{L{4.2cm} L{2.1cm} L{4.2cm} L{3.3cm}}
\toprule
\hdr{FACE-ATLAS claim} & \hdr{verdict} & \hdr{anchor support} & \hdr{honest tension}\\
\midrule
\rowcolor{facebg} Transdiagnostic dimensions over categories & supported & RDoC, HiTOP & ---\\
General factor read as \emph{burden}, not a symptom ``p'' & supported & Caspi p-factor (impairment) & p-factor's meaning is contested\\
\rowcolor{facebg} A graded \emph{continuum}, not discrete biotypes & open debate & normative models (Wolfers, Marquand) & B-SNIP biotypes (Clementz)\\
\textbf{Biology $\perp$ severity} & \textbf{calibrated} & transdiagnostic metabolic/immune; immuno-metabolic subtype & case-control elevations; biology tracks \emph{function}\\
\rowcolor{facebg} Metabolic and inflammatory are distinct axes & supported & different drivers / symptom maps & they co-cluster in a subgroup\\
\bottomrule
\end{tabular}
\end{table}

\subsection{The dimensional turn and the general factor}
The premise that severe mental illness is better described by continuous, diagnosis-crossing dimensions than by
DSM categories is now a mainstream research program, formalized in the NIMH Research Domain Criteria and the
Hierarchical Taxonomy of Psychopathology~\cite{insel2010rdoc,kotov2017hitop}. The general factor $G$ has a
direct precedent: the ``p factor'' of psychopathology, where a single general dimension fit the structure of
disorders better than internalizing/externalizing/thought-disorder factors and was read --- as $G$ is here ---
as \emph{life impairment} rather than a specific symptom liability~\cite{caspi2014pfactor}. That dimension is
also \emph{durable}: across four decades the diagnoses shift while the general dimension persists~\cite{caspi2020dunedin},
the same trait-vs-state pattern FACE-ATLAS reports longitudinally (Chapter~\ref{sec:temporal}). The bifactor
machinery itself is well-developed in psychometrics~\cite{reise2012bifactor,lahey2017general}.

The honest qualification is that the \emph{meaning} of a general psychopathology factor is contested:
symmetrical bifactor models can yield anomalous, sample-dependent solutions~\cite{heinrich2021bifactor}, and
multi-informant work suggests a single-rater general factor partly reflects method variance~\cite{watts2021pfactor}.
FACE-ATLAS sidesteps the worst of this by \emph{anchoring} $G$ on functioning and global-severity items only,
so it carries no symptom content (Chapter~\ref{sec:results}) --- a deliberately conservative reading that
answers the ``what does the general factor mean'' critique by construction rather than by interpretation.

\subsection{Continuum or biotypes? an unsettled debate}
FACE-ATLAS's conclusion that the transdiagnostic space is a graded continuum rather than a set of discrete
biotypes (Chapter~\ref{sec:strata}) enters a live and genuinely unresolved debate, and intellectual honesty
requires presenting the strongest opposing program as such. On the supporting side, individual-level
\emph{normative modelling} of brain structure in exactly the FACE disorders finds patient deviations so
heterogeneous that overlap exceeds a few percent at only a handful of loci --- ``the average patient'' is a
non-informative construct~\cite{wolfers2018normative}, a framing introduced as a deliberate alternative to
homogeneous case-control groups~\cite{marquand2016normative} and refined to show heterogeneity at the regional
scale but partial convergence at the network scale~\cite{segal2023heterogeneity}. On the opposing side, the
B-SNIP program used brain-based biomarkers to identify \emph{three replicable psychosis biotypes that ignore
DSM boundaries}~\cite{clementz2016biotypes} --- direct evidence for discrete biological subgroups.

\begin{caveatbox}[The biotype counter-evidence, stated fairly]
The B-SNIP biotypes are the most serious challenge to a pure-continuum reading, and we do not dismiss them.
Two points keep the accounts compatible rather than contradictory. First, they interrogate \emph{different
measurement spaces}: B-SNIP clusters cognitive/electrophysiological brain biomarkers, whereas FACE-ATLAS finds
no discrete structure on its nine clinical--biological dimensions (a five-method structure gate, adjusted Rand
index $\approx 0$ vs.\ DSM subtypes; Chapter~\ref{sec:strata}). Clusters in one space need not survive in the
other. Second, the B-SNIP data themselves described DSM \emph{diagnoses} as ``a single severity
continuum''~\cite{clementz2016biotypes}, and the program's recent synthesis now endorses \emph{both} categorical
and dimensional descriptions~\cite{clementz2024biotypes}. FACE-ATLAS is best read as a continuum claim about
\emph{its} feature space, consistent with the normative-modelling heterogeneity literature, not as a refutation
of brain-biomarker biotypes.
\end{caveatbox}

\subsection{The load-bearing claim: biology is largely severity-independent}
The headline of the measurement map --- that metabolic and inflammatory load are the least severity-entangled domains relative to the
general functional-burden factor ($\phi_{Gd}\approx 0.12$ and $0.07$; Chapter~\ref{sec:results}) --- is the
claim most in need of external corroboration, and the one we state most carefully. The published literature
does not report this exact contrast (a within-sample correlation between a biological-load factor and a
functional-burden factor is, to our reading, novel), but a large body of convergent evidence supports its
\emph{direction}, and a smaller body bounds its \emph{strength}. We therefore frame the finding as biology
being \textbf{largely independent of a general clinical/functional-burden axis}, not as strict statistical
orthogonality.

\textbf{Why the direction is well-grounded.} Cardiometabolic burden in severe mental illness is
\emph{transdiagnostic} and tracks age, body-mass index, medication and chronicity rather than diagnosis or
symptom severity: the largest meta-analyses find metabolic-syndrome and diabetes prevalence essentially equal
across schizophrenia, bipolar disorder and major depression, with multi-episode status and antipsychotic
exposure --- not severity --- the operative predictors~\cite{vancampfort2015metabolic,vancampfort2016diabetes}.
The biology is present \emph{before} chronic severity accrues: glucose dysregulation and insulin resistance
appear in antipsychotic-naïve, first-episode and even pre-onset cohorts~\cite{pillinger2017glucose,perry2019dysglycaemia,garridotorres2022metabolic}.
Inflammation, in turn, marks a \emph{subgroup} and \emph{specific symptoms} rather than the severity gradient:
only about a quarter of depressed patients are inflamed by a CRP$>$3\,mg/L threshold~\cite{osimo2019prevalence},
and where severity is tested directly, C-reactive protein tracks mood \emph{state} but not symptom severity in
bipolar disorder~\cite{fernandes2016bipolar} and positive --- not negative --- symptoms in
schizophrenia~\cite{fernandes2016schizophrenia}. Most strikingly, an independently developed framework ---
\emph{immuno-metabolic depression} --- holds that inflammation and metabolic abnormality co-cluster in a
$\sim$20--30\% subgroup mapping onto atypical, energy-related symptoms rather than overall
burden~\cite{penninx2024imd,milaneschi2020imd,lamers2013subtypes,lamers2020profilers}. That is, by entirely
different methods, the literature already describes biology as a partly separable dimension --- a strong
convergent validation of the FACE axis.

\textbf{Why ``independent'' must be qualified.} The tension is real and we state it. Case--control inflammatory
elevations are not null (CRP--depression $d\approx0.15$; larger in some
meta-analyses)~\cite{howren2009crp,osimo2020inflammation}; a modest severity/duration dose-response with
metabolic burden exists~\cite{penninx2018metabolic,pillinger2022variability}; and --- the most pointed caveat,
because $G$ is itself a \emph{functional}-burden axis --- metabolic dysfunction does track cognition and
functioning in mood disorders~\cite{maksyutynska2024cognition}. Genetic correlations between inflammation and
depressive symptoms are small but non-zero and largely carried by adiposity~\cite{kappelmann2021mr,khandaker2014il6},
which is why raw inflammation--severity associations attenuate after body-mass-index adjustment. A
FondaMental-adjacent study even reports a shared nutritional pathway loading onto \emph{both} symptoms and
metabolic comorbidity~\cite{faugere2025vitamins}, and treatment-emergent weight gain can co-occur with symptom
\emph{improvement}~\cite{luckhoff2018weight} --- a reminder that the biology--symptom relationship is
treatment-mediated, not a simple severity coupling.

\begin{keyresult}[The calibrated verdict on biology $\perp$ severity]
The external literature \emph{supports the direction} of the FACE finding --- metabolic and inflammatory load
are driven by medication, adiposity, chronicity and subgroup membership far more than by where a patient sits
on a global severity axis, and constitute a partly separable biological dimension (the immuno-metabolic-subtype
literature is almost a conceptual twin). It also \emph{bounds the strength}: case-control elevations and small,
adiposity-carried severity associations mean the relationship is weak and confounded, not zero. The defensible
claim is therefore \emph{``largely independent of a general clinical/functional-burden axis,''} with the
$\phi_{Gd}\approx 0.07$--$0.12$ values read as a \emph{quantification} of a pattern the prior literature had
described only qualitatively --- a contribution, and one the evidence base buttresses rather than contradicts.
\end{keyresult}

\subsection{What is novel here}
Situating the work also clarifies its contribution. Prior transdiagnostic-structure research is either
symptom-only (the p-factor and HiTOP traditions) or single-modality neuroimaging (the normative-model and
B-SNIP traditions); the immuno-metabolic literature asserts a separable biological dimension but, to our
reading, does not measure its correlation with a general bifactor. FACE-ATLAS instead places metabolic and
inflammatory laboratory markers \emph{directly into} the transdiagnostic factor space --- alongside cognition,
sleep, suicidality, mania and substance --- across three diagnostic cohorts at once, and \emph{quantifies} how
entangled each biological axis is with overall functional burden. That immune dysfunction is a large,
partly CNS-independent system in early psychosis~\cite{pillinger2019multisystem} makes such a measurement
worth having. The novelty claim is therefore narrow and defensible: FACE-ATLAS is, to our knowledge, the first
direct measurement --- with propagated uncertainty --- of how independent metabolic and inflammatory load are
from a transdiagnostic functional-burden axis, converting a qualitatively known pattern into a coordinate.
